\documentclass[../main.tex]{subfiles}\begin{document}

\section{Problemformulering}
Hvordan kan en digital platform understøtte og motivere lokale borgere til at hjælpe hinanden gennem både betalte hverdagsopgaver og frivillige aktiviteter, samtidig med at platformen skaber tillid mellem aktørerne og muliggør verificering og belønning af sociale bidrag?

\subsection{Underspørgsmål}
1. Hvordan kan en fælles platform understøtte forskellige typer af hjælp, herunder både betalte peer-to-peer-opgaver og frivillige aktiviteter?\\

2. Hvordan kan platformen skabe tillid mellem brugere, organisationer og virksomheder, når aktørerne ikke nødvendigvis kender hinanden på forhånd?\\

3. Hvordan kan udførelsen af frivillige aktiviteter verificeres på en måde, der reducerer risikoen for misbrug af platformens belønningssystem?\\

4. Hvordan kan et pointbaseret belønningssystem designes, så brugere belønnes for sociale bidrag, samtidig med at transaktioner og indløsninger håndteres konsistent og sporbar?\\

5. Hvordan kan platformen designes, så forskellige aktørtyper og fremtidige typer af aktiviteter og belønninger kan tilføjes uden væsentlige ændringer af systemets kerne?\\

\section{Gantt-diagram}
\begin{figure}[H]
    \centering
    \includegraphics[width=1\linewidth]{IMG/gantt_diagram.png}
    \caption{Gantt diagram for projektet}
    \label{fig:gantt}
\end{figure}

\section{Produktet + Aktører}

% Diagrams for clarity
\begin{figure}[H]
    \centering
    \includegraphics[width=1\linewidth]{IMG/actor.png}
    \caption{Aktører i produktet}
    \label{fig:actors}
\end{figure}

\section*{User Stories}

\subsection*{Epic 1 -- Konto \& Identitet}

\paragraph{US-01}
Som \textbf{user} vil jeg oprette en konto som \textbf{contributor/requester, organization} eller \textbf{reward partner}, så jeg kan deltage på platformen.

\paragraph{US-02}
Som \textbf{user} vil jeg vedligeholde en profil med relevante oplysninger om mig selv, så andre deltagere kan træffe en informeret beslutning om at interagere med mig.

\subsection*{Epic 2 -- Peer-to-peer-opgaver}

\paragraph{US-03}
Som \textbf{requester} vil jeg oprette en opgave, der beskriver den hjælp, jeg har brug for, så \textbf{contributors} kan tilbyde at udføre den.

\paragraph{US-04}
Som \textbf{contributor} vil jeg gennemse tilgængelige opgaver, så jeg kan finde opgaver, der er relevante for mig.

\paragraph{US-05}
Som \textbf{contributor} vil jeg afgive et tilbud på en opgave, så \textbf{requesteren} kan vælge mig til at udføre den.

\paragraph{US-06}
Som \textbf{requester} vil jeg vælge en \textbf{contributor} blandt de modtagne tilbud, så opgaven kan tildeles.

\paragraph{US-07}
Som \textbf{requester} og \textbf{contributor} vil jeg markere opgaven som fuldført, så transaktionen kan afsluttes.

\paragraph{US-08}
Som \textbf{user} vil jeg anmelde en anden deltager efter en fuldført opgave, så fremtidige brugere kan vurdere vedkommendes omdømme.

\paragraph{US-09}
Som \textbf{user} vil jeg se en anden brugers omdømme, så jeg kan vurdere, om jeg har tillid til vedkommende.

\subsection*{Epic 3 -- Aktiviteter med social impact}

\paragraph{US-10}
Som \textbf{organization} vil jeg oprette frivillige aktiviteter, så \textbf{contributors} kan finde måder at støtte vores sag på.

\paragraph{US-11}
Som \textbf{organization} vil jeg se, hvem der har tilmeldt sig min aktivitet, så jeg kan håndtere deltagelsen.

\paragraph{US-12}
Som \textbf{contributor} vil jeg gennemse tilgængelige frivillige aktiviteter, så jeg kan finde muligheder, der matcher mine interesser og min tilgængelighed.

\paragraph{US-13}
Som \textbf{contributor} vil jeg tilmelde mig en frivillig aktivitet, så \textbf{organizationen} ved, at jeg har til hensigt at deltage.

\paragraph{US-14}
Som \textbf{organization} vil jeg verificere, at en \textbf{contributor} har gennemført en aktivitet, så der kun tildeles Karma Koin for legitime bidrag.

\subsection*{Epic 4 -- Karma Koin}

\paragraph{US-15}
Som \textbf{contributor} vil jeg modtage Karma Koin efter en verificeret aktivitet, så mine sociale bidrag bliver anerkendt.

\paragraph{US-16}
Som \textbf{user} vil jeg se min nuværende Karma Koin-saldo, så jeg ved, hvor mange point jeg kan bruge.

\paragraph{US-17}
Som \textbf{user} vil jeg se min Karma Koin-transaktionshistorik, så jeg kan forstå, hvordan min saldo har ændret sig.

\paragraph{US-18}
Som \textbf{platform} vil jeg have, at enhver Karma Koin-transaktion kan spores tilbage til sin oprindelse, så belønningssystemet kan auditeres.

\subsection*{Epic 5 -- Belønninger}

\paragraph{US-19}
Som \textbf{reward partner} vil jeg oprette belønninger, så \textbf{contributors} kan indløse deres Karma Koin hos min virksomhed.

\paragraph{US-20}
Som \textbf{user} vil jeg gennemse tilgængelige belønninger, så jeg kan beslutte, hvordan jeg vil bruge mine Karma Koin.

\paragraph{US-21}
Som \textbf{user} vil jeg indløse Karma Koin til en belønning, så jeg får en fordel for mit bidrag.

\paragraph{US-22}
Som \textbf{reward partner} vil jeg se statistik for en given belønning, så jeg kan følge, hvor populære belønninger er sammenlignet med hinanden.

\subsection*{Epic 6 -- Tillid \& Sikkerhed}

\paragraph{US-23}
Som \textbf{platform} vil jeg have, at begrænsede handlinger kræver den rette rolle, så \textbf{users} ikke kan udføre handlinger, de ikke er autoriseret til.

\paragraph{US-24}
Som \textbf{administrator} vil jeg undersøge omstridte eller mistænkelige aktiviteter og transaktioner, så misbrug af platformen kan efterforskes.

\end{document}